\section{引言}

进入21世纪第三个十年，全球供应链体系经历了前所未有的系统性冲击。2020年初，新型冠状病毒疫情（COVID-19）导致全球制造业中心——中国——的大规模停工，引发从电子产品到医疗器械的全产业链断供危机，世界贸易组织估计2020年全球商品贸易量萎缩约5.3\%。随后，2022年俄乌冲突触发了能源与粮食供应链的地缘政治重构，欧洲天然气价格一度飙升十余倍。2023—2024年，巴拿马运河干旱与红海航运危机又叠加了气候与安全维度的物流中断。三波冲击的叠加效应，使得"供应链韧性"（Supply Chain Resilience）从运营管理的边缘议题跃升为各国经济安全议程的核心关切。

在这一宏观背景下，中国政府对供应链安全的重视程度迅速提升。2020年中央财经委员会第八次会议强调"构建现代流通体系"，2021年"十四五"规划首次将"产业链供应链韧性与安全水平"纳入国家发展目标，2023年《数字中国建设整体布局规划》进一步将数字化确立为赋能产业链供应链的核心手段。政策话语的密集转向意味着一个紧迫的实证问题亟待回答：企业数字化转型究竟能否提升供应链韧性？如果能，传导机制是什么？

学术界对这一问题的回应正处在从碎片化到系统化的过渡阶段。既有文献大致沿两条平行线推进。第一条线关注数字化对企业绩效的经济效应：大量研究证实了数字化投资对全要素生产率、创新产出和财务绩效的正向因果效应\cite{wu2025impact,chen2025supply}，但"供应链稳定性"作为结果变量长期被主流文献忽略。第二条线识别供应链韧性的决定因素，主要从供应商多元化、库存冗余、地理分散化和关系治理等运营策略出发\cite{ivanov2021digital,faruquee2021strategic}，数字化多被简化为一两个控制变量而非核心解释变量。真正将数字化作为供应链韧性核心驱动力的实证研究，在2024年以后才开始密集涌现\cite{li2025digital,zhang2024supply,liude2025supply}，且普遍存在一个关键局限：对传导机制的探讨多局限于单一中介路径（如供应链整合），缺乏对多条潜在机制的系统比较。对于数字化转型究竟通过何种渠道增强供应链韧性，以及哪条渠道的相对贡献更大，既有文献尚无法给出清晰的回答。

本文的边际贡献体现在以下三个方面。第一，在同一分析框架内同时提出并检验三条传导路径——流程优化度、企业创新能力和信息共享水平——系统地区分了"响应—适应—感知"三种韧性动态机制，首次为各路径的相对重要性提供了可比证据。第二，基于2015—2024年中国A股上市公司面板数据，采用企业—年份双向固定效应模型结合手动Sobel检验，为数字化转型与供应链韧性的因果关系提供了发展中国家的大规模微观证据。第三，结合COVID-19冲击这一外生事件，考察了数字化转型的供应链保护效应在系统性危机前后的时变特征，发现了一条具有政策含义的经验规律：数字化对韧性的提升效应在冲击期被显著放大。

本文其余部分安排如下：第二节为文献综述与理论假设，第三节介绍研究设计，第四节报告实证结果，第五节总结全文并提出政策建议。